\section{EPN: Event Proposal Network}

\subsection{Duration prediction}
EPN takes the three RED outputs per class and frame,
$[\hat p^{\pres}_{c,t},\hat p^{\on}_{c,t},\hat p^{\off}_{c,t}]$, and predicts
two non-negative durations:
\begin{align}
\hat d^{\on}_{c,t} &= \Softplus(h^{\on}_{c,t}),\\
\hat d^{\off}_{c,t} &= \Softplus(h^{\off}_{c,t}),
\end{align}
where $h$ is the output of a bidirectional GRU followed by a linear layer.
The proposal for frame $t$ is
\begin{equation}
\hat r_{c,t}
=
[t-\hat d^{\on}_{c,t},\; t+\hat d^{\off}_{c,t}].
\end{equation}
Two variants exist: \emph{per-class GRUs} (one GRU per class, input
$[0,1]^{3\times T}$) and a \emph{single GRU} over all classes (input
$[0,1]^{3C\times T}$). The per-class variant is used for AS-Strong-10; the
single-GRU variant is used for many classes. Both are implemented.

\subsection{IoU loss}
For a frame inside an event, the ground-truth interval is
\begin{equation}
r_{c,t}
=
[t-d^{\on}_{c,t},\; t+d^{\off}_{c,t}].
\end{equation}
Both $r_{c,t}$ and $\hat r_{c,t}$ contain $t$, so intersection and union have a
closed form:
\begin{equation}
\IoU_{c,t}
=
\frac{
\min(d^{\on}_{c,t},\hat d^{\on}_{c,t})
+
\min(d^{\off}_{c,t},\hat d^{\off}_{c,t})
}{
\max(d^{\on}_{c,t},\hat d^{\on}_{c,t})
+
\max(d^{\off}_{c,t},\hat d^{\off}_{c,t})
}.
\end{equation}
The loss is evaluated on active frames and weighted by ground-truth duration:
\begin{equation}
\mathcal{L}_{\mathrm{iou}}
=
\frac{1}{\sum_{c,t} y^{\pres}_{c,t}}
\sum_{c,t}
y^{\pres}_{c,t}
\left(1-\IoU_{c,t}\right)
\left(d^{\on}_{c,t}+d^{\off}_{c,t}\right).
\end{equation}
\begin{algorithm}[t]
\caption{Duration-weighted IoU loss for one batch}
\label{alg:iou}
\begin{algorithmic}[1]
\Require predicted $\hat d^{\on},\hat d^{\off}$; targets $d^{\on},d^{\off}$;
         presence labels $y^{\pres}$
\State $inter \gets \min(\hat d^{\on},d^{\on})+\min(\hat d^{\off},d^{\off})$
\State $union \gets \max(\hat d^{\on},d^{\on})+\max(\hat d^{\off},d^{\off})$
\State $iou \gets inter / \max(union,\epsilon)$ \Comment{$\epsilon=10^{-7}$}
\State $weight \gets d^{\on}+d^{\off}$
\State $L \gets \sum y^{\pres}(1-iou)\,weight
       \big/ \max(\sum y^{\pres},1)$
\State \Return $L$
\end{algorithmic}
\end{algorithm}

The implementation follows the algebra in Algorithm~\ref{alg:iou}. Both
intervals contain $t$, so the intersection and union reduce to sums of minima
and maxima; no interval sorting is needed. The duration factor weights every active frame by the length of its event, so
longer events contribute more to the sum than shorter ones. This is not an
equal-per-event weighting; it is the form used in the paper, where it was more
stable than direct duration regression. In our URBAN-SED runs the weight
mattered: $\lambda_{\mathrm{iou}}=2$ gave 25.87 PSDS1, while $0.5$ gave 35.42.
The paper tunes this value per backbone. On URBAN-SED the choice moves PSDS1 by
about 9.5 points, which is larger than the gain the method was meant to deliver;
$\lambda_{\mathrm{iou}}=4$ was not run, and no cross-dataset transfer claim is
made.

